\chapter*{How to Read This Book}
\addcontentsline{toc}{chapter}{How to Read This Book}
\markboth{How to Read This Book}{}

This handbook teaches competitive programming from the first loop to the techniques that decide olympiad
medals. It follows the USACO Guide syllabus (Bronze, Silver, Gold, Platinum, Advanced) and goes beyond it
wherever contests do. Every topic is presented the same way:

\begin{itemize}
  \item \textbf{What it is}, explained from scratch with a picture;
  \item \textbf{how to recognize} a problem that needs it (the part most books skip);
  \item the \textbf{key ideas, facts and complexity}, with proofs sketched where they build intuition;
  \item the \textbf{pitfalls} that cost people contests;
  \item \textbf{pro tricks}: less-known ideas that experienced competitors rely on;
  \item and the \textbf{template} in your library that implements it, with its keyboard shortcut.
\end{itemize}

\section*{Difficulty badges}

Each section carries a badge showing roughly where it sits:

\begin{center}
\begin{tabular}{ll}
\bronze & first steps: loops, simulation, brute force, simple math \\
\silver & sorting, binary search, prefix sums, basic graphs, greedy \\
\gold & DP, shortest paths, DSU, segment trees, number theory \\
\plat & advanced data structures, tree decompositions, DP optimizations \\
\ioi & beyond Platinum: flows, suffix structures, FFT, exotic tricks \\
\end{tabular}
\end{center}

If you are starting out, read Parts I and II in order and skim the rest to know what exists. If you are
experienced, jump to any chapter: they are self-contained, and the \emph{Pro trick} boxes are where most of
the non-obvious material lives. Part X, \emph{Practice}, solves seventeen classic problems step by step, from
reading the limits to tested code; try each one yourself before reading its solution.

\section*{Box legend}

\begin{idea}The central insight that makes the technique work.\end{idea}
\begin{recognize}Signals in a problem statement that point to this technique.\end{recognize}
\begin{pitfall}A mistake that produces wrong answers, TLE or runtime errors.\end{pitfall}
\begin{tip}Practical advice for writing or debugging the code.\end{tip}
\begin{trick}An advanced idea, often the difference between a good and a great solution.\end{trick}
\begin{fact}A useful number, identity or theorem to remember.\end{fact}
\usetemplate{Name of the template}{Ctrl+Alt+T, key}{cpprefix}

\section*{Conventions}

Arrays and vertices are 0-indexed; ranges are half-open, $[l, r)$ meaning $l \le i < r$, unless stated otherwise.
$n$ is the input size, $m$ the number of edges, $q$ the number of queries, $\log$ is base 2. Code is C++17 and
assumes your template (\cc{#define int long long}, \cc{ll}, \cc{vi}, \cc{all}, \ldots). The companion
\emph{User Guide} (GUIDE.pdf) explains the template, the library shortcuts and the tools.
